\paragraph{有限 jet 决定性、三模碰撞缺口、无限分支账本与 Smith/prime-axis 刚性}
在 fold-gauge 常数的 Bernoulli 递阶剥离、Fibonacci 共振双频的碰撞地板、自然扩张的分层 prime-register 外置、Smith 局部亏损谱的完全分类以及整数轴对齐椭圆的自由原子分解都已建立之后，还可进一步抽取出下列跨链结论。它们分别把有限奇指数 jet 的决定性、零模重整化的失效、无限深度历史的有限层账本不可能性、模 \(p^k\) 核增长谱对 Smith 正规形的完备判别，以及整数轴对齐椭圆幺半群上同态的 prime-axis 刚性统一为一组有限决定性命题。

\begin{theorem}[折叠规范常数的有限 \texorpdfstring{$\varphi$}{phi}-adic jet 完全决定有限偶 \texorpdfstring{$\zeta$}{zeta} 塔]\label{thm:xi-time-part63b-logcm-phiadic-jet-determines-even-zeta}
记
\[
q:=\frac{\varphi}{2}\in(0,1),
\qquad
L_m:=\log C_m-\log(2\pi).
\]
则对任意固定 \(R\ge 1\)，\(L_m\) 在精度
\[
O\!\bigl(q^{(2R+1)m}\bigr)
\]
下的截断渐近类唯一决定 \(B_2,\dots,B_{2R}\)，从而唯一决定
\[
\zeta(2),\zeta(4),\dots,\zeta(2R).
\]
等价地，fold-gauge 常数的有限 \(\varphi\)-adic jet 已经完整编码了前 \(R\) 个偶 Bernoulli 数与前 \(R\) 个偶值 \(\zeta\)-塔。
\end{theorem}
\begin{proof}
由定理 \ref{thm:fold-bin-gauge-constant-stirling-bernoulli-hierarchy}，对任意固定 \(R\ge 1\) 都有展开
\[
L_m
=
\sum_{r=1}^{R}
\frac{B_{2r}}{r(2r-1)}
\bigl(\varphi^{-1}+\varphi^{2r-3}\bigr)q^{(2r-1)m}
+O\!\bigl(q^{(2R+1)m}\bigr).
\]
因此该截断 jet 的自由系数恰为
\[
a_r:=
\frac{B_{2r}}{r(2r-1)}
\bigl(\varphi^{-1}+\varphi^{2r-3}\bigr)
\qquad(1\le r\le R).
\]
另一方面，定理 \ref{thm:xi-time-part60ab-gauge-constant-recursive-bernoulli-stripping} 已经给出逐阶极限消元反演：这些系数可由 \(L_m\) 递归无损恢复。又因
\[
\varphi^{-1}+\varphi^{2r-3}>0
\qquad(r\ge 1),
\]
故 \(a_r\) 与 \(B_{2r}\) 一一对应，从而前 \(R\) 个偶 Bernoulli 数由该 jet 唯一确定。

最后，由标准恒等式
\[
B_{2r}
=
(-1)^{r-1}\frac{2(2r)!}{(2\pi)^{2r}}\zeta(2r)
\qquad(r\ge 1),
\]
便可逐项恢复 \(\zeta(2),\dots,\zeta(2R)\)。证毕。
\end{proof}

\begin{corollary}[\texorpdfstring{$(\log C_m)$}{(log Cm)} 的完整渐近 germ 对偶值 \texorpdfstring{$\zeta$}{zeta} 塔是忠实不变量]\label{cor:xi-time-part63b-logcm-asymptotic-germ-faithful-even-zeta}
沿用定理 \ref{thm:xi-time-part63b-logcm-phiadic-jet-determines-even-zeta} 的记号。设 \((C_m)_{m\ge 1}\) 与 \((C_m')_{m\ge 1}\) 为两组折叠规范常数序列，并记
\[
L_m:=\log C_m-\log(2\pi),
\qquad
L_m':=\log C_m'-\log(2\pi).
\]
对每个 \(r\ge 1\)，记 \(B_{2r}(C)\)、\(\zeta_C(2r)\) 分别为由序列 \((C_m)\) 依定理 \ref{thm:xi-time-part63b-logcm-phiadic-jet-determines-even-zeta} 所恢复出的第 \(2r\) 个 Bernoulli 数与偶值 \(\zeta\)-数据；对 \((C_m')\) 亦作同样约定。
若对每个 \(R\ge 1\) 都有
\[
L_m-L_m'
=
O\!\bigl(q^{(2R+1)m}\bigr)
\qquad(m\to\infty),
\]
则
\[
B_{2r}(C)=B_{2r}(C'),
\qquad
\zeta_C(2r)=\zeta_{C'}(2r)
\qquad(\forall r\ge 1).
\]
换言之，\((\log C_m)\) 在全部奇指数层上的渐近 germ，对偶值 \(\zeta\)-塔给出一个忠实不变量。
\end{corollary}
\begin{proof}
任取 \(R\ge 1\)。假设 \(L_m\) 与 \(L_m'\) 在 \(O(q^{(2R+1)m})\) 精度上一致，则它们的前 \(R\) 阶 \(\varphi\)-adic jet 相同。由定理 \ref{thm:xi-time-part63b-logcm-phiadic-jet-determines-even-zeta}，两者的
\[
B_2,\dots,B_{2R}
\]
以及
\[
\zeta(2),\dots,\zeta(2R)
\]
必逐项一致。由于 \(R\) 任意，遂得对每个 \(r\ge 1\) 都有
\[
B_{2r}(C)=B_{2r}(C'),
\qquad
\zeta_C(2r)=\zeta_{C'}(2r).
\]
证毕。
\end{proof}

\begin{theorem}[三模不可去缺口与零模不可重整化]\label{thm:xi-time-part63b-three-mode-gap-zero-mode-nonrenormalizable}
设 \(\mathsf{Col}_m\) 为 bin-fold 的碰撞概率。则相对于均匀基线 \(F_{m+2}^{-1}\)，有
\[
\liminf_{m\to\infty}
F_{m+2}\left(\mathsf{Col}_m-\frac{1}{F_{m+2}}\right)
\ge
2C_\varphi^2
>
0.
\]
并且这一正下界已经由三模
\[
k=0,\qquad k=F_m,\qquad k=F_{m+1}
\]
单独强制。因而任何只依赖零模归一化的均匀化重整，都不可能消除此常数量级缺口。
\end{theorem}
\begin{proof}
由定理 \ref{thm:xi-time-part9p-chi2-ranktwo-resonant-floor}，
\[
\chi^2\!\bigl(p_m^{\mathrm{bin}}\Vert u_m\bigr)
=
F_{m+2}\mathsf{Col}_m-1,
\]
并且若删去两条 Fibonacci 共振频率 \(F_m,F_{m+1}\)，则仍有
\[
\chi^2\!\bigl(p_m^{\mathrm{bin}}\Vert u_m\bigr)-\chi^2_{\mathrm{cut}}(m)
=
a_m(F_m)^2+a_m(F_{m+1})^2
=
2C_\varphi^2+o(1).
\]
于是
\[
F_{m+2}\left(\mathsf{Col}_m-\frac{1}{F_{m+2}}\right)
=
\chi^2\!\bigl(p_m^{\mathrm{bin}}\Vert u_m\bigr)
\ge
2C_\varphi^2+o(1),
\]
取下极限即得所示不等式。

这里 \(k=0\) 模只贡献均匀基线
\(
F_{m+2}^{-1}
\),
而严格正的超额全部来自两条非零频率 \(F_m,F_{m+1}\)。因此，只在零模上作归一化而不消去这对 bulk 共振的任何方案，都不可能把碰撞概率重新压回
\(
F_{m+2}^{-1}+o(F_{m+2}^{-1})
\)
的尺度。证毕。
\end{proof}

\begin{theorem}[无限深度有限纤维塔的有限层 prime-register 不可能性]\label{thm:xi-time-part63b-infinite-depth-finite-prime-register-impossibility}
\leanverified{paper\_xi\_time\_part63b\_infinite\_depth\_finite\_prime\_register\_impossibility}
设
\[
X_0\xrightarrow{\pi_0}X_1\xrightarrow{\pi_1}X_2\xrightarrow{\pi_2}\cdots
\]
为可数集合上的有限纤维满射链，并假设对每个 \(j\ge 0\) 都存在 \(y_j\in X_{j+1}\) 使
\[
\#\pi_j^{-1}(y_j)\ge 2.
\]
记其逆极限为
\[
\widehat X
:=
\varprojlim(X_j,\pi_j)
=
\Bigl\{(x_0,x_1,x_2,\dots)\in \prod_{j\ge 0}X_j:\ \pi_j(x_j)=x_{j+1}\ \forall j\ge 0\Bigr\}.
\]
则在定理 \ref{thm:xi-layered-prime-register-sharp-natural-extension} 的分层 prime-register 外置语义下，不存在任何仅使用有限多个预先固定 prime slices 的精确单射实现，能够在 \(\widehat X\) 上同时保持全历史可恢复性。

更精确地，对任意 \(k\ge 1\)，深度 \(k\) 前缀的最坏情形精确恢复至少需要 \(k\) 个彼此独立的 prime slices；而在常系数情形 \(X_j=X\)、\(\pi_j=T\) 下，文稿中的自然扩张分支指标编码恰以一侧无穷序列实现这一极限最小方案。
\end{theorem}
\begin{proof}
固定 \(k\ge 1\)。把定理 \ref{thm:xi-layered-prime-register-sharp-natural-extension} 应用于截断链
\[
X_0\xrightarrow{\pi_0}X_1\xrightarrow{\pi_1}\cdots\xrightarrow{\pi_{k-1}}X_k.
\]
由于每一层都存在至少一个真实分支点，最坏情形下深度 \(k\) 的精确分层外置至少需要 \(k\) 个有效 prime slices；并且该定理同时给出恰在每层写入一枚有效素数的单射构造，说明这个下界是可达的。

反设整个逆极限 \(\widehat X\) 存在某个只使用有限多个 prime slices 的精确单射外置。把该编码限制到前 \(k\) 层坐标，便得到同一有限 prime-slice 预算下的深度 \(k\) 精确恢复方案。取 \(k\) 大于可用 prime slices 的总数，即与上面的必要性矛盾。故有限层 prime-register 不可能精确实现整个无穷历史。

在常系数情形 \(X_j=X\)、\(\pi_j=T\) 下，定理 \ref{thm:killo-natural-extension-branch-register} 已把自然扩张共轭为右移插入系统
\[
(x_0,x_{-1},x_{-2},\dots)\longmapsto
\bigl(x_0,(\beta(x_{-1}),\beta(x_{-2}),\dots)\bigr),
\]
每一步只写入一个新的分支指标；因此它正是有限深度最优外置在逆极限上的极限闭合。证毕。
\end{proof}

\begin{theorem}[Smith 等价的核增长谱判别准则]\label{thm:xi-time-part63b-smith-kernel-growth-spectrum-criterion}
\leanverified{paper\_xi\_time\_part63b\_smith\_kernel\_growth\_spectrum\_criterion}
设
\[
A,B\in M_{m\times n}(\ZZ),
\qquad
\operatorname{rank}_{\QQ}(A)=\operatorname{rank}_{\QQ}(B)=m.
\]
则下列条件等价：
\[
\#\ker(A\bmod p^k)=\#\ker(B\bmod p^k)
\qquad
(\forall p\ \text{素数},\ \forall k\ge 1);
\]
\[
A,\ B\ \text{具有完全相同的 Smith 不变量}.
\]
换言之，整数线性投影在全部模 \(p^k\) 上的核增长谱，恰是其 Smith 正规形的完备判别式。
\end{theorem}
\begin{proof}
设 \(A\) 的 Smith 不变量为
\[
d_1^{A}\mid \cdots \mid d_m^{A},
\]
并记
\[
e_i^{A}(p):=v_p(d_i^{A}).
\]
对每个素数 \(p\) 与整数 \(k\ge 1\)，由定理 \ref{thm:xi-kernel-loss-divisibility-valuation-rational-euler}，
\[
\#\ker(A\bmod p^k)
=
p^{\,k(n-m)+\sum_{i=1}^{m}\min(k,e_i^{A}(p))}.
\]
因此若定义
\[
f_{A,p}(k):=
\log_p\#\ker(A\bmod p^k)-k(n-m),
\]
则
\[
f_{A,p}(k)=\sum_{i=1}^{m}\min(k,e_i^{A}(p)).
\]
于是其一阶差分满足
\[
\Delta_{A,p}(k):=
f_{A,p}(k)-f_{A,p}(k-1)
=
\#\{i:\ e_i^{A}(p)\ge k\}.
\]
故序列 \(\bigl(\#\ker(A\bmod p^k)\bigr)_{k\ge 1}\) 唯一决定多重集
\[
\{e_i^{A}(p)\}_{i=1}^{m}
\]
；逐素数拼合便唯一决定 \(A\) 的全部 Smith 不变量。对 \(B\) 亦然。

因此，若 \(A\) 与 \(B\) 的核增长谱逐点相同，则对每个 \(p\) 都有
\[
f_{A,p}(k)=f_{B,p}(k)
\qquad(\forall k\ge 1),
\]
从而
\[
\{e_i^{A}(p)\}_{i=1}^{m}
=
\{e_i^{B}(p)\}_{i=1}^{m}.
\]
逐素数合并便得两者 Smith 正规形相同。反向 implication 显然：相同的 Smith 不变量给出相同的模 \(p^k\) 核大小公式。证毕。
\end{proof}

\begin{theorem}[信息亏损 \texorpdfstring{$\zeta$}{zeta} 函数的 Euler 因子化与 Smith 刚性]\label{thm:xi-time-part63b-loss-zeta-euler-factorization-smith-rigidity}
设
\[
A\in M_{m\times n}(\ZZ),
\qquad
\operatorname{rank}_{\QQ}(A)=m,
\]
其 Smith 不变量为
\[
d_1\mid\cdots\mid d_m.
\]
对每个素数 \(p\) 记
\[
e_i(p):=v_p(d_i).
\]
定义信息亏损 \(\zeta\)-函数
\[
\mathcal Z_A(s)
:=
\sum_{b\ge 1}\frac{\#\ker(A_b)}{b^{\,s+n-m}}
\qquad(\Re s\gg 1).
\]
则有 Euler 因子化
\[
\mathcal Z_A(s)
=
\prod_p
\left(
\sum_{k\ge 0}
p^{-ks+\sum_{i=1}^{m}\min(k,e_i(p))}
\right).
\]
特别地，对同型尺寸的满秩整数矩阵 \(A,B\)，有
\[
\mathcal Z_A=\mathcal Z_B
\quad\Longleftrightarrow\quad
A,\ B\ \text{具有相同的 Smith 正规形}.
\]
\end{theorem}
\begin{proof}
由推论 \ref{cor:killo-kernel-size-general-modulus}，
\[
\#\ker(A_b)
=
b^{\,n-m}\prod_{i=1}^{m}\gcd(d_i,b).
\]
因此归一化系数
\[
a_A(b):=
\frac{\#\ker(A_b)}{b^{\,n-m}}
=
\prod_{i=1}^{m}\gcd(d_i,b)
\]
是乘法函数，从而 \(\mathcal Z_A(s)=\sum_{b\ge 1}a_A(b)b^{-s}\) 具有 Euler 乘积。对 \(b=p^k\) 进一步有
\[
a_A(p^k)
=
\prod_{i=1}^{m}p^{\min(k,e_i(p))}
=
p^{\sum_{i=1}^{m}\min(k,e_i(p))},
\]
故逐素数展开正是
\[
\mathcal Z_A(s)
=
\prod_p
\left(
\sum_{k\ge 0}
p^{-ks+\sum_{i=1}^{m}\min(k,e_i(p))}
\right).
\]

若 \(\mathcal Z_A=\mathcal Z_B\)，则两者的 Dirichlet 系数逐项相同，特别地对每个素数幂 \(p^k\) 都有
\[
\#\ker(A\bmod p^k)\,p^{-k(n-m)}
=
\#\ker(B\bmod p^k)\,p^{-k(n-m)}.
\]
于是
\[
\#\ker(A\bmod p^k)=\#\ker(B\bmod p^k)
\qquad(\forall p,\ \forall k\ge 1).
\]
由定理 \ref{thm:xi-time-part63b-smith-kernel-growth-spectrum-criterion}，\(A\) 与 \(B\) 具有相同的 Smith 正规形。反向 implication 显然：相同的 Smith 不变量给出相同的 Euler 因子。证毕。
\end{proof}

\begin{proposition}[整数轴对齐椭圆单子的 prime-axis 刚性]\label{prop:xi-time-part63b-integer-axis-ellipse-prime-axis-rigidity}
沿用定理 \ref{thm:xi-integer-ellipse-free-commutative-monoid} 的记号，令
\[
\mathsf E_{\NN}
:=
\{E_{a,b}:x^2/a^2+y^2/b^2=1,\ a,b\in\NN_{>0}\},
\qquad
E_{a,b}\odot E_{c,d}:=E_{ac,bd}.
\]
则
\[
(\mathsf E_{\NN},\odot)
\cong
\bigl(\NN^{(\mathcal P)}\times \NN^{(\mathcal P)},+\bigr).
\]
因而对任意交换幺半群 \((M,+)\)，每个幺半群同态
\[
J:(\mathsf E_{\NN},\odot)\to (M,+)
\]
都唯一由素数轴生成元
\[
E_{p,1},\qquad E_{1,p}
\qquad(p\in\mathcal P)
\]
的像所决定。特别地，当 \(M=\RR\) 时，存在唯一系数族 \((\alpha_p)_{p\in\mathcal P}\)、\((\beta_p)_{p\in\mathcal P}\) 使得对任意 \(a,b\in\NN_{>0}\) 都有
\[
J(E_{a,b})
=
\sum_{p}\alpha_p\,v_p(a)
+
\sum_{p}\beta_p\,v_p(b).
\]
\end{proposition}
\begin{proof}
定理 \ref{thm:xi-integer-ellipse-free-commutative-monoid} 已给出幺半群同构
\[
(\NN^{(\mathcal P)}\times \NN^{(\mathcal P)},+)
\cong
(\mathsf E_{\NN},\odot),
\qquad
(r_1,r_2)\longmapsto E_{N(r_1),N(r_2)}.
\]
其中 \(\NN^{(\mathcal P)}\) 是以素数集 \(\mathcal P\) 为基的自由交换幺半群。故 \((\mathsf E_{\NN},\odot)\) 本身就是以
\[
\{E_{p,1},E_{1,p}:p\in\mathcal P\}
\]
为原子集的自由交换幺半群。

自由对象上的幺半群同态由生成元像唯一决定，于是任意
\[
J:(\mathsf E_{\NN},\odot)\to (M,+)
\]
都完全由 \(J(E_{p,1})\) 与 \(J(E_{1,p})\) 唯一决定。

当 \(M=\RR\) 时，记
\[
\alpha_p:=J(E_{p,1}),
\qquad
\beta_p:=J(E_{1,p}).
\]
又由唯一原子分解
\[
E_{a,b}
=
\left(\bigodot_{p}E_{p,1}^{\odot v_p(a)}\right)
\odot
\left(\bigodot_{p}E_{1,p}^{\odot v_p(b)}\right),
\]
应用幺半群同态性即可得到
\[
J(E_{a,b})
=
\sum_{p}\alpha_p\,v_p(a)
+
\sum_{p}\beta_p\,v_p(b).
\]
由于原子分解唯一，系数族亦唯一。证毕。
\end{proof}

\endinput
